\section{Introduction} \label{sec:introduction}
% NLU, applications: challenges and concerns
Natural language understanding (NLU) is a core module in advanced artificial intelligence language processing systems such as sentiment analysis \cite{socher2013recursive} for customer feedback, intent classification \cite{wang2018glue} for conversational agents, semantic similarity \cite{cer2017semeval} for information retrieval, and so on. According to the modern approach, these sytems have achieved impressive capabilities by leveraging large language models \cite{devlin2019bert}. However, they often require centralized training and deployment, which poses challenges related to data privacy, computational cost, and inference latency.

% SLMs: advantages, challenges and concerns
Consequently, recent research has shifted towards small language models (SLMs) \cite{sanh2019distilbert, turc2019well, wang2020minilm} due to their low inference latency and cost-effectiveness. These models are well-suited to resource-constrained environments, addressing the challenges of large language models \cite{li2020federated}. Nevertheless, SLMs are often limited in capacity, making it difficult to maintain high performance across a variety of NLU tasks.

%MTL: advantages, challenges and concerns
At the same time, multi-task learning (MTL) \cite{ruder2017overview, caruana1997multitask} has been widely adopted as an effective paradigm to solve multiple tasks simultaneously. This approach of leveraging useful information from related tasks can lead to better learning performance \cite{wang2018glue, yu2020gradient}. Despite its advantages, most existing multi-tasking learning methods typically require centralized training data, which is impractical in privacy-sensitive situations where data is distributed across multiple clients.

% FL: advantages, challenges and concerns
With the growing availability of sensitive textual data, strict privacy regulations such as the GDPR \cite{european2016regulation} have raised concerns about centralizing user data in conventional machine learning pipelines. These challenges have motivated the development of federated learning (FL), a distributed paradigm that enables model training across clients without transferring raw data to a central server, thereby preserving privacy while supporting large-scale learning \cite{mcmahan2017communication, yang2019federated}. 
However, current FL pipelines mainly depend on batch synchronization where clients execute multiple local steps and will execute the sync periodically \cite{mcmahan2017communication, li2020federated_prox}. This design is effective for static datasets but less so for streaming and dynamic environments where language patterns as well as subjects and user intent change with frequency. 

% Our proposal
To address the aforementioned challenges and concerns, this paper proposes RT-FedMTL, a novel real-time federated MTL framework with SLMs for NLU tasks. Particularly, our framework integrates FL and MTL to enable collaborative training across distributed clients while preserving data locality and leveraging shared representations across tasks. To the best of our knowledge, this is the first work to explore the use of WebSocket protocols and SLMs tuning to reduce the feedback loop and provide near-instantaneous global updates for a NLU-based federated MTL system. Our key contributions are
summarized as follows:
\begin{enumerate}
    \item We propose RT-FedMTL, a novel real-time federated framework of MTL and SLMs to harness cross-task knowledge transfer without causing a large amount of resources to be consumed on devices.
    \item We assess five SLM architectures (DistilBERT \cite{sanh2019distilbert}, BERT-Medium \cite{turc2019well}, BERT-Mini \cite{turc2019well}, MiniLM \cite{wang2020minilm}, and TinyBERT \cite{jiao2019tinybert}) with respect to three standard NLU tasks (sentiment analysis, paraphrase detection, and semantic similarity), to get the best trade-off between performance and resource usage.
    \item We investigate the effects of non-IID and task heterogeneity on convergence speed and task accuracy in real-time domain.
    \item We show RT-FedMTL can drastically minimize per-client GPU memory requirements (up to 74-93\%) versus centralized MTL baselines, while keeping performance for multiple benchmarks (GLUE suite).
\end{enumerate}

The remainder of the paper is organised as follows. Section~\ref{sec:problem_definition} defines the formal problem of RT-FedMTL and Section~\ref{sec:proposed_method} presents the proposed framework. Section~\ref{sec:experiments} describes the experimental setup and results analysis. The relevant work is briefly reviewed in Section~\ref{sec:related_work}. 
Finally, Section~\ref{sec:conclusion} concludes and outlines future work.
